% ============================================================
% Rebar in Concrete: How Process Knowledge Binds a
% Cross-Platform Cognitive Network
%
% Paper 4 in the Cognitive MRI series
% ============================================================

\documentclass[11pt,twocolumn]{article}

% --- Page geometry ---
\usepackage[margin=0.75in]{geometry}

% --- Math ---
\usepackage{amsmath,amssymb}

% --- Tables ---
\usepackage{booktabs}
\usepackage{multirow}

% --- Graphics ---
\usepackage{graphicx}

% --- Typography ---
\usepackage{microtype}

% --- Compact lists ---
\usepackage{enumitem}
\setlist{nosep,leftmargin=*}

% --- Citations ---
\usepackage{natbib}
\bibliographystyle{plainnat}

% --- Cross-references and links ---
\usepackage[breaklinks=true]{hyperref}
\hypersetup{
  colorlinks=true,
  linkcolor=blue!70!black,
  citecolor=blue!70!black,
  urlcolor=blue!70!black,
}
\usepackage{xurl}

% --- Colors for figures ---
\usepackage{xcolor}
\definecolor{cgblue}{HTML}{4A90D9}
\definecolor{agred}{HTML}{D94A4A}

% --- Captions ---
\usepackage[font=small,labelfont=bf]{caption}

% --- Compact title ---
\usepackage{titling}
\setlength{\droptitle}{-2em}
\pretitle{\begin{center}\Large\bfseries}
\posttitle{\end{center}\vspace{-1em}}
\preauthor{\begin{center}\normalsize}
\postauthor{\end{center}\vspace{-1em}}
\predate{\begin{center}\small}
\postdate{\end{center}\vspace{-1.5em}}

\begin{document}

% ============================================================
% Title block
% ============================================================
\title{Rebar in Concrete: How Process Knowledge Binds\\a Cross-Platform Cognitive Network}

\author{
  Alexander Towell\thanks{Corresponding author. Email: atowell@siue.edu. ORCID: 0000-0001-6443-9897}
  \quad
  John Matta \\
  \small Department of Computer Science,
  Southern Illinois University Edwardsville
}

\date{}

\maketitle

% ============================================================
\begin{abstract}
% ============================================================
When the same person explores knowledge through both conversational AI (ChatGPT) and agentic AI (Claude Code) over two years, do the resulting cognitive traces converge or remain platform-specific?
We construct a concept-level network from 115 cognitive concepts extracted via LLM abstraction from 1,908 ChatGPT conversations and 3,945 Claude Code sessions belonging to a single user.
Concepts are embedded in a shared 768-dimensional space and connected by cosine similarity ($\theta = 0.70$), yielding a small-world network ($\sigma = 2.42$, $\omega = 0.03$) with 1,280 edges.
We find that \emph{community membership carries $48.5\times$ more information about a concept's neighbors than platform identity does} (normalized MI $= 0.162$ vs.\ $0.007$).
Platform origin is virtually invisible to network structure.
Percolation analysis reveals an asymmetric architecture: as the similarity threshold rises, agentic process concepts are the last to disconnect ($81.8\%$ of the residual giant component at $\theta = 0.85$), yet the network survives complete agentic removal with $66.1\%$ of its giant component intact.
Like rebar in concrete, process knowledge forms the strongest internal bonds while domain knowledge provides the structural mass.
Cross-platform links are the network's most robust feature: zero of nine local bridges and zero of eight true bridges are cross-platform, and cross-platform edge fraction remains stable at $\sim$$40\%$ across all thresholds (CV $= 0.062$).
Topological features alone predict cross-platform edges at AUC $= 0.956$.
These findings demonstrate that a single user's cognitive structure, viewed through two AI modalities, forms a unified conceptual space organized by meaning rather than platform.
\end{abstract}

% ============================================================
\section{Introduction}
\label{sec:intro}
% ============================================================

The past three years have seen an explosion in AI-assisted knowledge work.
Millions of users maintain conversation archives with ChatGPT, Claude, and other large language models~\citep{zhao2024wildchat}, creating externalized records of their cognitive activity~\citep{zhao2023survey}.
More recently, \emph{agentic} AI tools---Claude Code, GitHub Copilot, Cursor---have extended this interaction beyond conversation: these systems read files, write code, execute tests, and spawn autonomous subagents, leaving behind session logs that document not just what was discussed but what was \emph{done}~\citep{wang2024survey}.

In prior work, we introduced a ``cognitive MRI'' methodology that transforms AI conversation archives into semantic similarity networks~\citep{towell2025cognitive,towell2026temporal}.
Applied to a single user's ChatGPT history, this revealed 15 knowledge communities, super-linear densification ($\gamma = 1.405$), and stable community structure over 29 months.
The session-level (``episodic'') networks, however, face a fundamental limitation: they can only compare sessions \emph{within} a single platform.
A ChatGPT conversation about Bayesian inference and a Claude Code session implementing a Bayesian model live in separate networks with no structural connection.

This paper introduces a \emph{concept layer} that bridges this gap.
From the episodic communities detected in each platform's similarity network, we extract abstract cognitive concepts via LLM-based abstraction---labels like ``Bayesian inference as model selection'' or ``Iterative Refinement.''
This approach is methodologically related to text network analysis tools like InfraNodus~\citep{paranyushkin2019infranodus}, which maps discourse structure via word co-occurrence networks, and to the community-level LLM summarization pipeline of GraphRAG~\citep{edge2024graphrag}.
These concepts are then embedded in a shared semantic space and connected by cosine similarity, producing a single cross-platform network where every concept competes on equal footing regardless of its platform of origin.

The resulting network reveals a striking finding: \textbf{platform identity is virtually invisible to network structure.}
Knowing which community a concept belongs to provides $48.5\times$ more information about its neighbors than knowing which platform it came from.
Agentic process concepts (iteration, debugging, refactoring) and ChatGPT knowledge concepts (Bayesian reasoning, statistical modeling, philosophical inquiry) intermix freely, organized by semantic affinity rather than platform provenance.

Yet the two platforms contribute differently.
As we raise the similarity threshold and the network fragments, \emph{agentic concepts are the last to disconnect}---they form the strongest-bonded core.
But the network survives the complete removal of all agentic concepts with most of its giant component intact, because ChatGPT knowledge concepts provide the structural mass.
This asymmetry evokes a material science metaphor: agentic process knowledge is the \emph{rebar} in the \emph{concrete} of domain knowledge---the strongest internal material, binding the structure, but embedded within a bulk that provides the volume and the load-bearing capacity.

Our contributions are:
\begin{enumerate}
  \item A methodology for constructing cross-platform concept networks from heterogeneous AI archives.
  \item Quantitative evidence that a single user's cognitive structure transcends platform boundaries (platform NMI $= 0.007$).
  \item Characterization of an asymmetric ``rebar in concrete'' architecture where process and domain knowledge play structurally complementary roles.
  \item Demonstration that cross-platform integration is the network's most robust structural feature, with zero fragile cross-platform links and topologically predictable edges (AUC $= 0.956$).
\end{enumerate}


% ============================================================
\section{Methods}
\label{sec:methods}
% ============================================================

\subsection{Episodic layer construction}

The source data comprise two archives from a single user: 1,908 ChatGPT conversations spanning December 2022 to April 2025 (29 months) and 3,945 Claude Code sessions spanning November 2025 to March 2026 (101 days, of which 1,061 are human-initiated parent sessions and 2,884 are AI-spawned subagents).

Following the methodology of~\citet{towell2025cognitive}, we embed each session using nomic-embed-text~\citep{nussbaum2024nomic,reimers2019sentence} (768 dimensions), compute all-pairs cosine similarity, and threshold to construct episodic similarity networks.
ChatGPT conversations use user:AI message weighting of 2:1 and threshold $\theta = 0.90$ (validated by ablation study); Claude Code sessions use text-only content mode and threshold $\theta = 0.95$ (adjusted for the higher baseline similarity of coding sessions).
These episodic networks serve as the substrate from which concepts are extracted.

Community detection via the Louvain algorithm~\citep{blondel2008fast,newman2006modularity,fortunato2010community} identifies 15 ChatGPT communities (e.g., ML/AI, Statistics, Philosophy, Web Development) and 12 agentic communities (mapped to coding projects rather than knowledge domains).

\subsection{Concept extraction}
\label{sec:methods:concepts}

From each episodic community, we extract abstract cognitive concepts using LLM-based abstraction (llama3.2 via Ollama).
For each community, the top 20 sessions by message count are presented to the LLM with a prompt requesting 2--5 abstract concepts that characterize the community's cognitive content---an approach analogous to topic modeling via latent Dirichlet allocation~\citep{blei2003latent}, but leveraging LLM comprehension rather than bag-of-words statistics.
This yields 79 ChatGPT concepts and 36 agentic concepts (115 total).

ChatGPT concepts are predominantly \emph{knowledge-domain} concepts: ``Bayesian Reasoning for Model Selection,'' ``Statistical Modeling and Estimation,'' ``Philosophical Inquiry.''
Agentic concepts are predominantly \emph{process-workflow} concepts: ``Iterative Refining,'' ``Debugging Iterative Refinement,'' ``Verification-Driven Development.''

Stability validation (five extractions per community, three LLM models) confirms that concept \emph{themes} are robust---Bayesian reasoning always emerges from the statistics community---though exact phrasing varies.
Semantic deduplication at $\theta \geq 0.85$ recovers 2.7 stable concepts per community.

\subsection{Concept network construction}
\label{sec:methods:network}

Each concept is represented as ``name: description'' and embedded via nomic-embed-text into the same 768-dimensional space used for episodes~\citep{veremyev2019graph}.
All $\binom{115}{2} = 6{,}555$ pairwise cosine similarities are computed.
Edges are placed between concepts with similarity exceeding a threshold $\theta$.

We select $\theta = 0.70$ as the operating point based on a systematic sweep (Figure~\ref{fig:robustness}):
at this threshold the network has meaningful modularity ($Q = 0.276$), a large giant component ($96\%$), and substantial cross-platform community mixing ($57\%$ of communities contain both ChatGPT and agentic concepts).
Lower thresholds produce trivial community structure ($Q < 0.07$); higher thresholds fragment the network excessively.

The resulting concept network has 115 nodes, 1,280 edges, 7 Louvain communities, and density 0.195.
Table~\ref{tab:network_summary} summarizes key properties.

\begin{table}[t]
  \centering
  \caption{Concept network summary ($\theta = 0.70$).}
  \label{tab:network_summary}
  \begin{tabular}{lr}
    \toprule
    Property & Value \\
    \midrule
    Nodes (ChatGPT / Agentic) & 115 (79 / 36) \\
    Edges & 1,280 \\
    Density & 0.195 \\
    Giant component & 95.7\% \\
    Communities (Louvain) & 7 \\
    Modularity $Q$ & 0.276 \\
    Small-world $\sigma$ & 2.42 \\
    Small-world $\omega$ & 0.03 \\
    Avg.\ path length & 2.07 \\
    Clustering coefficient & 0.598 \\
    Diameter & 7 \\
    Cross-platform edge fraction & 39.1\% \\
    \bottomrule
  \end{tabular}
\end{table}

\subsection{Analysis methods}

We employ the following analytical tools:
\begin{itemize}
  \item \textbf{Information theory}: Mutual information between platform/community labels and neighbor labels, normalized by marginal entropy~\citep{danon2005comparing}.
  \item \textbf{Percolation analysis}: Threshold sweep from $\theta = 0.45$ to $0.95$, tracking giant component fraction and platform composition~\citep{stauffer1994introduction,albert2000error}.
  \item \textbf{Dimensionality reduction}: PCA and t-SNE~\citep{vandermaaten2008visualizing} on the similarity matrix.
  \item \textbf{Structural vulnerability}: Local bridge detection (edges with no common neighbors), true bridge detection (edges whose removal disconnects the graph).
  \item \textbf{Link prediction}: Logistic regression on topological features (Adamic-Adar~\citep{adamic2003friends,libennowell2003link}, common neighbors, Jaccard index, degree sum)~\citep{lu2011link}.
  \item \textbf{Core-periphery}: $k$-core decomposition and Borgatti-Everett continuous model~\citep{borgatti2000models}.
  \item \textbf{Rich-club analysis}: Normalized rich-club coefficient $\rho(k)$ against 100 configuration model realizations~\citep{colizza2006detecting}.
  \item \textbf{Null models}: Erd\H{o}s-R\'enyi~\citep{erdos1960evolution} (same density), configuration model (same degree sequence), geometric random graph~\citep{penrose2003random} (same edge count)---100 realizations each.
  \item \textbf{Communicability}: Matrix-exponential communicability and Estrada index~\citep{estrada2008communicability}.
\end{itemize}


% ============================================================
\section{Results}
\label{sec:results}
% ============================================================

We organize results around three questions: (1) How does the concept network encode platform identity? (2) What structural role do agentic process concepts play? (3) How robust is cross-platform integration?

\subsection{Platform invisibility}
\label{sec:results:invisible}

The concept network's most striking property is the near-complete absence of platform signal in its structure.

\paragraph{Information content.}
Mutual information analysis reveals that community membership carries $48.5\times$ more information about a concept's neighbor identity than platform does (Table~\ref{tab:information}).
The normalized MI for platform is $0.007$---observing a concept's neighborhood tells you essentially nothing about whether it came from ChatGPT or Claude Code.
In contrast, community NMI ($0.162$) provides meaningful predictive power, confirming that the network is organized by semantic content, not origin.

\begin{table}[t]
  \centering
  \caption{Mutual information between node attributes and neighbor attributes.  Community membership is $48.5\times$ more informative than platform.}
  \label{tab:information}
  \begin{tabular}{lccc}
    \toprule
    Attribute & MI (bits) & $H$ (bits) & Norm.\ MI \\
    \midrule
    Platform   & 0.007 & 0.900 & 0.007 \\
    Community  & 0.316 & 1.949 & 0.162 \\
    \midrule
    Ratio (community / platform) & \multicolumn{3}{c}{$\mathbf{48.5\times}$} \\
    \bottomrule
  \end{tabular}
\end{table}

\paragraph{Geometric embedding.}
Dimensionality reduction confirms this finding geometrically.
In t-SNE projection of the concept similarity matrix, $69.4\%$ of agentic concepts fall \emph{inside} the convex hull of ChatGPT concepts (Figure~\ref{fig:tsne}).
Communities separate $4.5\times$ more cleanly than platforms (mean intra-class $-$ inter-class similarity: $0.057$ for community vs.\ $0.013$ for platform).
Agentic concepts are literally embedded within the ChatGPT knowledge cloud, not forming a separate region.

\paragraph{Communicability.}
The ratio of cross-platform to same-platform communicability (matrix exponential) is $1.005$---information traverses the network with effectively identical efficiency regardless of platform boundaries~\citep{estrada2008communicability}.
The concept network is a single communicability basin.

\begin{figure}[t]
  \centering
  \fbox{\parbox{0.9\linewidth}{\centering\small\textit{Figure placeholder: t-SNE projection of 115 concepts.
  ChatGPT concepts in {\color{cgblue}blue}, agentic in {\color{agred}red}.
  Dashed line: ChatGPT convex hull enclosing 69.4\% of agentic points.
  Communities shown as background shading.}}}
  \caption{t-SNE embedding of the concept space.  Agentic process concepts (red) are geometrically embedded within the ChatGPT knowledge cloud (blue).  Community structure (shaded regions) provides $4.5\times$ stronger separation than platform identity.}
  \label{fig:tsne}
\end{figure}


\subsection{The rebar architecture}
\label{sec:results:rebar}

Although platform is invisible to the \emph{static} network, the two platforms contribute asymmetrically to its structural backbone.
We characterize this asymmetry through percolation, resilience, and triadic analysis.

\paragraph{Percolation dominance.}
As the similarity threshold rises above $\theta = 0.70$, the network fragments.
Figure~\ref{fig:percolation} tracks the giant component fraction and its agentic composition.
At the critical threshold $\theta_c \approx 0.80$, the giant component drops below $50\%$.
At $\theta = 0.85$, only 11 concepts remain connected---and $81.8\%$ of them are agentic (9 agentic, 2 ChatGPT), compared to the baseline composition of $31.3\%$.
The strongest concept-level connections in the entire network are between agentic process concepts: debugging links to testing links to refactoring with similarities above $0.85$.

\paragraph{Resilience asymmetry.}
Despite forming the strongest core, agentic concepts are not structurally necessary.
Removing \emph{all 36 agentic concepts} leaves the giant component at $66.1\%$ of the original---the ChatGPT knowledge backbone remains self-sufficient.
Conversely, it requires removing $70\%$ of ChatGPT concepts to halve the giant component.

This is the rebar-in-concrete asymmetry: agentic process concepts are the strongest-bonded material, binding distant knowledge domains through shared cognitive patterns (iteration, decomposition, abstraction).
But ChatGPT knowledge concepts provide the structural mass---the volume and interconnection that hold the network together even without the rebar.

\begin{figure}[t]
  \centering
  \fbox{\parbox{0.9\linewidth}{\centering\small\textit{Figure placeholder: Percolation cascade.
  Left axis: GC fraction (solid line) vs.\ threshold.
  Right axis: agentic \% in GC (dashed line).
  Dramatic crossover at $\theta = 0.80$: GC drops below 50\% while agentic \% rises to 81.8\% at $\theta = 0.85$.}}}
  \caption{Percolation analysis of the concept network.  As threshold rises, the giant component fragments---but agentic concepts are the last to disconnect.  At $\theta = 0.85$, the residual GC is $81.8\%$ agentic.}
  \label{fig:percolation}
\end{figure}

\paragraph{Triadic structure.}
Triangle census~\citep{milo2002network} provides a motif-level view of the rebar function.
Compared to degree-preserving configuration models, the concept network has $3.0\times$ more triangles than expected ($z = 48.4$).
Critically, all four platform-composition types (CCC, CCA, CAA, AAA) are equally over-represented ($z = 20$--$28$), confirming that triangle formation is topological, not a platform-segregation artifact.

At the per-node level, agentic concepts participate in cross-platform triangles at a dramatically higher rate: $67\%$ of agentic triangles involve at least one ChatGPT concept, versus $42\%$ for ChatGPT concepts ($p < 0.0001$).
Agentic concepts rarely close triads with only other agentic concepts; they exist to stitch ChatGPT knowledge into cross-platform structures.

\subsection{Cross-platform mirrors}
\label{sec:results:mirrors}

We identify cross-platform ``mirror'' pairs: the best 1-to-1 matching between each ChatGPT concept and its most similar agentic counterpart.
Of 36 possible pairs, 14 are ``strong mirrors'' with similarity $\geq 0.80$ (Table~\ref{tab:mirrors}).

The strongest mirror is ``Bayesian inference as model selection,'' independently extracted from both platforms with near-identical embedding (sim $= 0.943$)---the second-strongest edge in the \emph{entire} network.
Other strong mirrors include ``Problem Decomposition'' ($0.908$), ``Knowledge Graphs / Knowledge Graphing'' ($0.916$), and ``Pattern of Abstraction / Abstraction and Simplification'' ($0.877$).
These are genuine cognitive primitives: the same reasoning patterns encoded identically regardless of whether the user encountered them through conversational exploration or engineering practice.

Strong mirror pairs occupy the network's most central positions: $1.9\times$ higher degree, $1.9\times$ higher betweenness, and $1.5\times$ higher $k$-core than unmatched concepts ($p < 10^{-6}$).
Though they represent only $24\%$ of nodes ($14$ pairs $= 28$ concepts), they form the structural backbone.

\begin{table}[t]
  \centering
  \caption{Top cross-platform concept mirrors.  All five strongest mirrors occupy the same Louvain community.}
  \label{tab:mirrors}
  \begin{tabular}{p{2.4cm}p{2.4cm}r}
    \toprule
    ChatGPT concept & Agentic concept & Sim \\
    \midrule
    Bayesian inference as model selection & Bayesian inference as model selection & .943 \\
    Knowledge Graphs \& Network Analysis & Knowledge Graphing & .916 \\
    Problem Decomposition & Problem Decomposition & .908 \\
    Iterative Dev.\ \& Refining & Iterative Problem-Solving & .894 \\
    Pattern of Abstraction & Abstraction \& Simplification & .877 \\
    \bottomrule
  \end{tabular}
\end{table}


\subsection{Structural robustness}
\label{sec:results:robust}

Cross-platform integration is not merely present---it is the network's most robust structural feature, by several independent measures.

\paragraph{Zero fragile cross-platform links.}
Of 9 local bridges (edges sharing no common neighbors) and 8 true bridges (edges whose removal disconnects the graph), \emph{none} are cross-platform.
All structurally fragile edges are same-platform peripheral connections linking niche concepts within ChatGPT or within the agentic corpus.
Cross-platform edges are deeply embedded: they share on average $36.8\%$ of their endpoints' neighbors, slightly more than same-platform edges ($35.2\%$, $p = 0.066$).

\paragraph{Threshold invariance.}
Sweeping $\theta$ from $0.55$ to $0.85$, the cross-platform edge fraction remains remarkably stable: $35$--$44\%$ with coefficient of variation $0.062$ (Figure~\ref{fig:robustness}).
This is the most robust metric in the entire analysis; by comparison, community count varies with CV $= 1.24$ and modularity with CV $= 0.73$.
At the extreme tail ($\theta \geq 0.90$), cross-platform pairs reach $50\%$, because the very highest similarities are identical cognitive primitives appearing on both platforms.

\begin{figure}[t]
  \centering
  \fbox{\parbox{0.9\linewidth}{\centering\small\textit{Figure placeholder: Cross-platform fraction vs.\ threshold.
  Horizontal band at $\sim$40\% with CV $= 0.062$.
  Comparison: modularity Q (CV $= 0.73$) and community count (CV $= 1.24$) overlaid as dashed lines showing high variability.}}}
  \caption{Threshold robustness of cross-platform edge fraction.  The $\sim$$40\%$ cross-platform fraction is the most stable network metric, barely varying from $\theta = 0.55$ to $0.85$.}
  \label{fig:robustness}
\end{figure}

\paragraph{Topological predictability.}
Logistic regression predicts cross-platform edge existence from topological features alone at AUC $= 0.956$ (Table~\ref{tab:prediction}).
The dominant predictor is the Adamic-Adar index~\citep{adamic2003friends}---concept pairs sharing many well-connected common neighbors are virtually guaranteed to be linked, regardless of platform.
A single feature (Adamic-Adar) achieves AUC $= 0.951$, nearly matching the full eight-feature model.
Cross-platform links are structurally \emph{necessary}: they connect concepts that should be connected based on their neighborhood topology.

\begin{table}[t]
  \centering
  \caption{Cross-platform edge prediction from topology.  A single topological feature achieves AUC $= 0.951$.}
  \label{tab:prediction}
  \begin{tabular}{lr}
    \toprule
    Model & AUC \\
    \midrule
    Full model (8 features) & $0.956 \pm 0.015$ \\
    Adamic-Adar alone & 0.951 \\
    Common neighbors alone & 0.946 \\
    Degree sum alone & 0.868 \\
    \bottomrule
  \end{tabular}
\end{table}


\subsection{Core-periphery organization}
\label{sec:results:core}

The concept network exhibits clear core-periphery structure with a cross-platform gradient.

\paragraph{$k$-core decomposition.}
Using $k$-core decomposition as a measure of structural importance~\citep{kitsak2010identification}, we find the innermost core ($k = 19$, 31 nodes, density $0.82$) contains 19 ChatGPT and 12 agentic concepts---a $61$:$39$ ratio, compared to the full network's $69$:$31$.
Agentic concepts are disproportionately represented in the core.
Cross-platform edge fraction increases from $39.1\%$ (full network) to $47.3\%$ (inner cores), confirming that the core is where platforms converge.

\paragraph{Borgatti-Everett coreness.}
Continuous core-periphery modeling~\citep{borgatti2000models} assigns coreness scores $c_i \in [0, 1]$.
The top-scoring concepts are ``Iterative Refining'' (agentic, $c = 0.925$) and ``Problem-Solving Frameworks'' (ChatGPT, $c = 0.911$).
Coreness correlates with cross-platform degree at $\rho = 0.894$ ($p < 10^{-40}$)---being structurally central is nearly synonymous with being cross-platform connected.

\paragraph{Rich-club effect.}
High-degree concepts interconnect $20$--$34\%$ more densely than expected from their degree sequence (normalized $\rho = 1.2$--$1.3$, $z > 20$; \citealt{colizza2006detecting}).
Within this rich club, agentic concepts are more cross-platform (8 of 10 connections cross platforms) than ChatGPT concepts (3 of 10), quantifying the rebar role at the elite level.

\paragraph{Flat topology.}
The concept network contains \emph{zero} hub nodes (all within-module $z$-scores $< 2.5$ in the Guimer\`a-Amaral classification; \citealt{guimera2005functional}).
Bridging is distributed across 37 satellite connectors---filling structural holes~\citep{burt1992structural}---rather than concentrated in a few hubs, yielding a leaderless, egalitarian architecture.


\subsection{Null model validation}
\label{sec:results:null}

Three null model batteries confirm the significance of the observed structure.

\begin{itemize}
  \item \textbf{Clustering}: $3.0\times$ more triangles than the configuration model ($z = 48.4$); $2.8\times$ more than Erd\H{o}s-R\'enyi ($z = 60.1$).  The clustering is genuine, not an artifact of the degree sequence.
  \item \textbf{Cross-platform mixing}: The observed $39.1\%$ is \emph{lower} than the configuration model null ($43.5\%$, $p < 0.001$), meaning same-platform concepts are slightly more similar than cross-platform concepts.  The null model correction is important: cross-platform integration is real but is not \emph{stronger} than chance---it is structured differently from chance, with cross-platform edges concentrated in the core rather than uniformly distributed.
  \item \textbf{Network class}: Clustering falls between Erd\H{o}s-R\'enyi ($z = +60$) and geometric random graphs ($z = -9.4$), placing the concept network in the ``small-world semantic space'' regime characteristic of cognitive and semantic networks~\citep{steyvers2005large,dedeyne2019small}---more structured than random but less constrained than geometry.
\end{itemize}


% ============================================================
\section{Discussion}
\label{sec:discussion}
% ============================================================

\subsection{A unified cognitive space}

The central finding of this study is that a single user's cognitive concepts, extracted from two different AI platforms over two years, form a \emph{unified} network organized by semantic content rather than platform origin.
Platform identity carries only $0.7\%$ of the information that community membership carries about network neighborhoods.

This finding connects to a rich tradition of cognitive network science~\citep{siew2019cognitive}.
\citet{steyvers2005large} showed that semantic networks derived from population-level data (WordNet, free associations) exhibit small-world structure with power-law degree distributions; \citet{kenett2014investigating} demonstrated that individual differences in semantic network topology correlate with creative ability.
Our concept network extends this program to a novel domain: concepts extracted not from controlled experiments but from naturalistic AI-assisted knowledge work.
The small-world signature we observe ($\sigma = 2.42$) is consistent with the semantic network class identified by Steyvers and Tenenbaum, suggesting that AI-mediated concept formation follows the same organizational principles as human semantic memory.

This extends the ``extended mind'' thesis~\citep{clark1998extended,hutchins1995cognition}: not only do external AI tools function as cognitive extensions, but the \emph{conceptual structures} they help construct converge across tools.
The user does not maintain separate ``ChatGPT knowledge'' and ``Claude Code knowledge''---they maintain a single cognitive repertoire that manifests through whatever tool is at hand.

This convergence is not trivially expected.
ChatGPT concepts are predominantly \emph{epistemic} (``Bayesian inference,'' ``Statistical Modeling''), while agentic concepts are predominantly \emph{practical} (``Iterative Refining,'' ``Debugging'').
That these semantically different concept types intermix into a single network---with $40\%$ of edges crossing the platform boundary at every threshold---suggests that the underlying cognitive operations (abstraction, decomposition, iteration) transcend the knowledge/practice distinction.

\subsection{The rebar metaphor}

The asymmetric architecture we describe differs from the ``hub-and-spoke'' or ``rich-get-richer'' models common in network science~\citep{barabasi1999emergence}.
There are no hubs.
Instead, the two platforms contribute complementary structural materials:

\begin{itemize}
  \item \textbf{Concrete (ChatGPT knowledge)}: 79 concepts providing the bulk and the load-bearing structure. The network survives complete agentic removal. Knowledge concepts form the mass of the giant component at every threshold.
  \item \textbf{Rebar (agentic process)}: 36 concepts forming the strongest internal bonds. They dominate the residual core under percolation ($81.8\%$ at $\theta = 0.85$) and participate in cross-platform triangles at $67\%$ vs.\ $42\%$. They bind the concrete together.
\end{itemize}

This metaphor has a precise structural interpretation.
In reinforced concrete, the rebar's tensile strength resists cracking while the concrete's compressive strength bears the load.
In the concept network, process concepts' high mutual similarity ($> 0.85$) resists fragmentation under threshold pressure, while knowledge concepts' broad coverage maintains connectivity.
Neither material functions alone: unreinforced concrete cracks (a knowledge-only network has no strong cross-domain links); rebar without concrete collapses (the 36 agentic concepts alone form a tiny, disconnected graph).

\subsection{Structural inevitability of cross-platform links}

The AUC of $0.956$ for cross-platform edge prediction from topology alone (Section~\ref{sec:results:robust}) is perhaps the most consequential result for network theory.
It means that cross-platform edges are not random additions to the network---they are \emph{topologically required}.
The Adamic-Adar predictor's dominance indicates that cross-platform links arise where two concepts share many well-connected common neighbors.
In other words, the network generates its own cross-platform integration through local structure.

This has implications for the study of multiplex networks more broadly~\citep{boccaletti2014structure,kivela2014multilayer,battiston2014structural}.
When different layers of a system share an underlying semantic space, inter-layer links may be structurally determined rather than requiring a separate generative mechanism.
The concept layer does not need a ``cross-platform attachment rule''---it needs only cosine similarity in a shared embedding space, and the cross-platform structure emerges from the geometry.

\subsection{Comparison with multiplex network literature}

The ``single communicability basin'' finding (cross/same ratio $= 1.005$) contrasts with multiplex systems where inter-layer communication is substantially weaker than intra-layer.
In brain networks, structural and functional layers show complementary but not identical organization~\citep{bullmore2009complex}.
In social games, positive, negative, and economic layers encode largely distinct structure~\citep{szell2010multirelational}.
Our concept network is more integrated than these systems: once concepts are embedded in a shared space, platform identity genuinely vanishes from the communication dynamics.

However, the null model correction (Section~\ref{sec:results:null}) tempers this conclusion.
Cross-platform mixing at $39\%$ is \emph{below} the $43.5\%$ configuration model null, meaning same-platform concepts are slightly more similar than cross-platform concepts.
The integration is real but does not \emph{exceed} what the degree sequence alone would produce; rather, it is \emph{structured differently}, concentrated in the core ($47\%$) rather than uniformly distributed.

\subsection{Limitations}

\textbf{Single subject.}
This study analyzes one user's AI archives.
Whether the rebar architecture generalizes to other users---or reflects this individual's particular cognitive style---remains open.

\textbf{LLM-mediated concept extraction.}
Concepts are extracted by a language model (llama3.2), introducing a layer of interpretation.
Exact concept names vary across runs (10\% exact-match stability), though semantic themes are robust.
Different extraction LLMs produce varying concept inventories for technical domains, though philosophical/abstract concepts show near-perfect cross-model agreement ($> 0.95$).

\textbf{Threshold dependence.}
While the cross-platform fraction is stable (CV $= 0.062$), community structure varies substantially with threshold (CV $= 1.24$).
The specific community assignments should be interpreted with caution.

\textbf{Asymmetric source corpora.}
The ChatGPT archive spans 29 months with 1,908 conversations; the Claude Code archive spans 3.4 months with 3,945 sessions.
The $79$:$36$ concept ratio reflects this asymmetry.
A longer agentic observation period might alter the balance.

\textbf{No causal claims.}
The convergence we observe is correlational.
We cannot determine whether the user's ChatGPT experience \emph{caused} their Claude Code concepts to align, whether both reflect pre-existing cognitive structure, or whether the shared embedding model imposes convergence.


% ============================================================
\section{Conclusion}
\label{sec:conclusion}
% ============================================================

We have shown that cognitive concepts extracted from a single user's two-year engagement with both conversational and agentic AI form a unified semantic network in which platform identity is virtually invisible ($48.5\times$ less informative than community membership).
The network's architecture is asymmetric: agentic process concepts form the strongest internal bonds (dominating the residual core under percolation) while ChatGPT knowledge concepts provide the structural mass (the network survives complete agentic removal).
Cross-platform integration is the network's most robust feature: zero fragile cross-platform links, stable $\sim$$40\%$ fraction across all thresholds, and near-perfect topological predictability (AUC $= 0.956$).

The rebar-in-concrete metaphor captures a structural principle that may extend beyond this case study: when a person uses multiple cognitive tools, \emph{process knowledge} (how to think) may bind \emph{domain knowledge} (what to think about) into a coherent whole.
The concepts that transcend platforms---iteration, decomposition, abstraction, pattern recognition---are not specialized to any domain or tool.
They are the cognitive rebar: the strongest, most cross-cutting, most structurally central elements of the knowledge network.

Future work should test this architecture across multiple users, longer time horizons, and additional AI platforms.
If the rebar-in-concrete pattern generalizes, it would suggest that AI-assisted cognition converges on platform-independent conceptual structures---that the tools we use to think become invisible once we extract the thoughts themselves.


% ============================================================
% References
% ============================================================
\bibliography{rebar-refs}

\end{document}
